\section{Scalar order twenty-four and a smaller exceptional set}

These changes retain the global arithmetic hypotheses and conventions of
the pinned Version 10 source, commit
\texttt{036cebe4f1a4aa7ca89820a9b48f3629b4c364d6}. They concern the
scalar valuation argument, the set of excluded rational primes, and the
numerical anchor. They do not independently verify the inherited geometric,
modularity, height, prime-pattern, group, or pole theorems.

\subsection{Replace the scalar order thirty-two by twenty-four}

Retain
\[
 \beta=2401+2754(2X+1)^2,\qquad
 r=\beta(\beta-1)(32\beta+343)>0,\qquad c=42.
\]
Let $J=H=[h]P$ and $Q=S-P=[a-1]P$, both nonidentity. The scalar
norm tree uses the two rational coefficients of $f(J),f(Q)$ at order
one, and those of $L=z(Q)/z(J)$ at a larger order $m$.

The proof in \texttt{manuscripts/quartic\_v10.tex}, lines
10422--10548, is uniform in these positive integer orders. At a consumed
prime let $R=v_p(r)$ be odd. The pole obstruction and the constant leaf
on the all-left path force $R>0$ and $v_p(W)>-R/2$.
For a norm node $N_r(u,v)=u^2+rv^2$, the two nonzero summands have
opposite valuation parity. The implication
\[
 v_p(N_r(u,v))>-tR,\quad t<1
 \quad\Longrightarrow\quad
 v_p(u)>-tR/2,\quad v_p(v)>-(t+1)R/2
\]
therefore descends through the fixed finite tree. Every normalized
leaf $X_i/r^{m_i}$ has valuation greater than $-R$, so
\[
 v_p(X_i)>(m_i-1)R\ge m_i-1,
 \qquad v_p(X_i)\ge m_i.
\]
The first four leaves put $J,Q$ in the formal group at every place above
$p$. Formal logarithm isometry and $[h]Q=[a-1]J$ identify the valuation
of $L$ with that of $(a-1)/h$. Thus the two ratio leaves of order $m$
give $v_p(a-1)\ge m$, and
\[
 M_5(\beta)^m\mid a-1,\qquad
 h(\beta)\le\mathcal H |a-1|^{19/m}.
\]
Here the radical and height constant may be either the old ones or the
smaller-set ones established below.

Take $m=24$. Put $q=r^{-1}$, already present for the order-one leaves,
and impose
\[
 rq_3=q^2,\qquad q_6=q_3^2,\qquad
 q_{12}=q_6^2,\qquad q_{24}=q_{12}^2.
\]
Since $r>0$, these uniquely give $q_j=r^{-j}$. There are four new
coordinates instead of the five squarings needed for $r^{-32}$.
The direct ratio numerator is multiplied by $q_{24}$ instead of
$r^{-32}$. All equations remain quadratic, including when the anchor
coefficient is lifted to a coordinate.

The final soundness proof needs $h(\beta)<|a|/40$. Write
$C=|a-1|$. Its existing anchor margin is $C>(80\mathcal H)^{10}$.
Since $1-19/24=5/24>1/10$,
\[
 \mathcal H C^{19/24}<C/80\le |a|/40.
\]
Thus the same Lucas and input-height contradictions apply without
changing their detector orders. In particular the input comparison may
use order twelve: the printed bound is
\[
 \log B<\frac{12}{5}\log(40\mathcal H)
             +\frac{19}{5}h(\beta)
 <\frac{31}{200}|a|<|a|-3.
\]
The Lucas residual depends only on $a,t,b$, whose comparisons retain
order forty.

Completeness changes only finite lower valuation bounds for the scalar
leaves. The old order-thirty-two preparation already supplies the weaker
order-twenty-four bounds. At the finite support of $42r$, the new norm
tree remains a unit by its same constant all-left leaf. The inherited
pole-completeness and continuity arguments therefore apply. This is a
weaker scalar predicate with a checked global consumer, not a claimed
bijection between the old and new pole witness spaces.

An isolated rebuild gives 422 total variables, 390 quadratic residuals,
and 26,603 quartic monomials. Order thirty-one also saves one variable,
using four Laurent quotients, but has 26,639 monomials. Order twenty-four
was selected because it gives the same variable saving with fewer terms.

\subsection{Use the finite exceptional set actually needed}

Define
\[
 \mathcal S=\{p:p\le41\}\cup\{109,439,2399\},
 \qquad B_{\mathcal S}=\prod_{p\in\mathcal S}p
 =34926271784737221178290.
\]
Let $M_{\mathcal S}(t)$ be the product of primes outside $\mathcal S$
having odd contact with one of
$\infty,0,1,2401,-343/32$. The proposed replacement height constant is
\[
 A_{19}=2^{109}210^{1800},\qquad
 \mathcal H_{\mathcal S}=2A_{19}B_{\mathcal S}^{19}+128.
\]
The exact integer has bit length 15,419. The simpler bound sufficient
below is
\[
 \mathcal H_{\mathcal S}<2^{15936},
\]
because $B_{\mathcal S}<2^{75}$ and
$2+109+1800\cdot8+19\cdot75=15936$.

This does not follow merely by changing the symbol $B_*$ in a theorem
statement. It follows by reapplying its actual proof, as follows.

The integral geometry in lines 6935--7003 is smooth away from $210$.
The five branch sections are disjoint there; their determinant primes
lie in $\{2,3,5,7\}$. Tame local coordinates show that even contact
introduces no ramification after an unramified extension. Consequently,
for the same normal fiber field and surface,
\[
 \operatorname{rad}(D_LN_{\rm bad})
       \mid B_{\mathcal S}M_{\mathcal S}(t).
\]
The arguments (LR1)--(LR5), lines 7090--7405, use only a squarefree
$R$ containing this ramification support and the fixed primes dividing
$210$. Their character-order, field-degree and conductor bounds are
therefore unchanged with $R=B_{\mathcal S}M_{\mathcal S}(t)$:
\[
 d'\le2304R,\qquad N\le210^{300}R^2,
 \qquad\operatorname{rad}(N)\mid R.
\]
The exponent-nineteen proof in lines 8915--9073 then gives the same
$h_F(A)<A_{19}R^{19}$ and geometric comparison. Hence, subject to the
same inherited global inputs,
\[
 h(t)\le\mathcal H_{\mathcal S}M_{\mathcal S}(t)^{19}
 \qquad(t\notin\{0,1,2401,-343/32\}).
\]

\paragraph{Actual local consumers.}
The following verifications justify consuming the larger radical.

The pole obstruction excludes negative odd valuations at every prime
outside $\{2,3,17\}$ (lines 7698--7718). The one-place norm argument
(lines 11032--11060) only additionally needs the unramifiedness of
$\mathbb Q(\sqrt2)$ and integral normalized valuations. These hypotheses
hold outside $\mathcal S$; that proof has no further numerical lower
bound on the prime.

For the detector contact map $\phi(t)=2+1/t$, the actual contact and
guard-difference primes are explicitly listed in lines 11103--11113:
\[
 \{2,3,5,7,11,37,109,439,2399\}\subset\mathcal S.
\]
All fixed coefficients needed to identify the detector radicals and
transfer the three scale heights are units away from this set.
The proof of the product-height estimate in
\texttt{docs/quartic-523-540/addition.tex}, lines 284--349, uses the same
three scales and is unchanged with $\mathcal H_{\mathcal S}$.

The new scalar product has only the three contacts $0,1,-343/32$.
Their collisions involve $2,3,5,7$, and $2754=2\cdot3^4\cdot17$ is a
unit outside $\mathcal S$. The square fiber at 2401 and the even
negative valuations of $\beta$ therefore prove the same odd-contact
implication for $M_{\mathcal S}(\beta)$. The historical extra factor
$2\beta-1$, with its collision primes, is absent from this source.
The anchor excludes $X=-1/2$ before this height theorem is used.

Only the actual logarithm orders twelve and forty are needed. The
formal identity and coefficient bound in lines 9220--9258 give
\[
 \log_E Q=zG(w),\quad v(w)=4r,\quad
 v(a_j)\ge-\frac{4j+1}{p-1},\quad r=v(z)\ge1.
\]
For $m\in\{3,10\}$, consider $j\ge m$. If $4j+1<p$, the coefficient
is integral and the term has valuation at least $4mr$. Otherwise,
$p>41$ implies $j\ge11$ and $p-1\ge42$. The inequality
\[
 \frac{4j+1}{42}\le4(j-10)\le4(j-m)r
\]
then gives the same valuation lower bound. These bounds tend to infinity,
so they control the entire tail. Thus $G_2$ has relative precision twelve
and $G_9$ has relative precision forty at every prime outside
$\mathcal S$, including inert primes. The same proof with $m=1$
gives $\log_E(Q)/z(Q)\in1+p^{4r}\mathcal O$.
The prime 41 is deliberately excluded: the coefficient $a_{10}$ has
41-adic valuation $-1$, so a depth-one point would only give order 39
for that first omitted term.

All two-anchor comparisons now use these specialized estimates and
the fact that one of $h,h-1$ is a unit. No blanket strengthening of the
historical order-52 or order-64 formal-logarithm statements is asserted.
For a primary identity input the earlier zero-ratio argument still
applies at the appropriate order twelve or forty.

\paragraph{Completeness.}
Keep the existing prime-pattern choice of contact primes larger than
76176, including all its finite exclusions and Gaussian congruences.
Shrinking the exceptional set in soundness does not require choosing
smaller primes in completeness. The formulas express each detector's
odd valuations through exactly its chosen prime contact values; elsewhere
the valuations are even, apart from fixed exceptional primes already in
$\mathcal S$. Thus no new uncontrolled small odd contact appears.
The prime-separation, progression, small-prime pole and real estimates
are unchanged. The scalar leaf order is the finite value twenty-four,
so its completeness and continuity arguments above apply as well.

\subsection{Shorten the anchor and its constant chains}

Set
\[
 C_{\rm a}=3^{2^{18}},\qquad S_0=C_{\rm a}^2=3^{2^{19}}.
\]
Retain the same field-norm physical anchor and quaternion conic, changing
only this fixed coefficient. The local argument is uniform in its
exponent: it gives the same odd-root and local-input conditions, while
the scalar shifted point obeys
\[
 0\ne a-1\in2\,3^{2^{17}-1}\mathbb Z.
\]
The exact subgroup calculation is in lines 8132--8209 of the source;
it says $v_3(f([s]P))=2+2v_3(s)$ for nonzero even $s$.
Consequently
\[
 \log_2|a-1|\ge1+(2^{17}-1)\log_2 3
 >196607.5>159430
 =10(7+15936)>\log_2(80\mathcal H_{\mathcal S})^{10}.
\]
This supplies every numerical margin used by the unchanged Lucas
contradiction and the order-twelve input contradiction. The scalar
order-twenty-four bound supplies $h(\beta)<|a|/40$ as already shown.

The standard decimal chain now starts at $3^{2^{10}}$ and ends at
$3^{2^{19}}$, using ten coordinates instead of seventeen. The larger-
coefficient chain starts at $3^{2^{18}}$ and uses two coordinates instead
of nine. The fixed-coefficient core count is unchanged by the anchor
reduction itself. The producer, specialization verifier, modular-content
checker, lift construction and lift checker must all use the declared
new exponent; changing only a printed parameter name is insufficient.

The provided tools parse that exponent from the IR metadata. Their sign
certificates use $S_0>2^{2^{19}}$ in this instance, and the constant chains
are checked against their actual endpoint. The isolated scalar-order-
twenty-four plus anchor-nineteen branch has 422 core coordinates,
432 standard decimal coordinates, and 424 larger-coefficient coordinates.
